\subsection{Contrastive Weight as Speed Control}
\label{sec:disc-whisper}

A natural question arises from the four-regime taxonomy: does the
contrastive weight $c_w$ control \emph{whether} block-diagonal structure
emerges, or \emph{how quickly} it emerges? The FI-004 annealing experiment
(\S\ref{sec:disc-annealing}) found an optimal fixed
$c_w \approx 0.02$ and a cliff at $c_w = 0$, suggesting a minimum
threshold. But FI-004 used fixed weights---the Steersman's adaptive decay
was not tested in the low-$c_w$ regime.

CW-001 tests this directly. Configuration: $n = 6$ (RD), identity
initialization, initial $c_w = 0.02$ with the Steersman's adaptive
decay active. Control Law~2 reduces $c_w$ when the co/cross ratio
exceeds the directionality threshold, with a floor at
$0.25 \times c_{w,\text{base}} = 0.005$. The experiment runs for
10{,}000~steps on TinyLlama~1.1B.

\subsubsection{Three-Phase Trajectory}

The CW-001 trajectory reveals three distinct dynamical regimes within
a single training run:

\begin{table}[h]
\centering
\small
\begin{tabular}{@{}rrrrl@{}}
\toprule
Step & Co/Cross & Fiedler & $c_w$ & Phase \\
\midrule
100 & 20:1 & 0.0073 & 0.020 & 1: Initial growth \\
500 & 1{,}592:1 & 0.0004 & 0.015 & 1: Rapid BD formation \\
1{,}500 & 1{,}867:1 & 0.0016 & 0.005 & 2: Plateau onset \\
3{,}300 & 1{,}855:1 & 0.0018 & 0.005 & 2: Sustained plateau \\
4{,}600 & 2{,}252:1 & 0.0019 & 0.005 & 3: Breakout \\
5{,}800 & 3{,}860:1 & 0.0014 & 0.005 & 3: Acceleration \\
7{,}500 & 7{,}884:1 & 0.0007 & 0.005 & 3: Exponential growth \\
8{,}300 & \textbf{15{,}183:1} & \textbf{0.0004} & 0.005 & 3: Peak \\
9{,}000 & 13{,}793:1 & 0.0004 & 0.005 & 4: BD stabilization \\
9{,}500 & 12{,}846:1 & 0.0004 & 0.005 & 4: Oscillation \\
\textbf{10{,}000} & \textbf{13{,}456:1} & \textbf{0.0005} & 0.005 & \textbf{4: Final} \\
\bottomrule
\end{tabular}
\caption{CW-001 trajectory: co/cross ratio, Fiedler eigenvalue, and
contrastive weight across four dynamical phases.}
\label{tab:cw001-trajectory}
\end{table}

\textbf{Phase~1 (steps 0--1{,}500): Rapid initial growth.}
The Steersman decays $c_w$ from 0.020 to the floor of 0.005 within
1{,}500~steps as the co/cross ratio rises rapidly. BD structure forms
quickly during this high-$c_w$ window, reaching ${\sim}1{,}800$:1.

\textbf{Phase~2 (steps 1{,}500--4{,}200): Apparent plateau.}
With $c_w$ floored at 0.005, the co/cross ratio oscillates between
1{,}500:1 and 2{,}000:1 for nearly 3{,}000~steps. The Fiedler
eigenvalue rebounds from its Phase~1 minimum, rising to 0.0020---indicating
that spectral connectivity is \emph{increasing} even as BD strength
appears stalled. This phase was initially interpreted as a terminal
state (adaptive decay defeating whisper-strength). It is not.

\textbf{Phase~3 (steps 4{,}600--8{,}300): Exponential breakout.}
The co/cross ratio breaks above 2{,}000:1 at step~4{,}600 and enters an
exponential growth phase: $2{,}252 \to 3{,}860 \to 7{,}884 \to 12{,}921
\to 15{,}183$:1 (peak at step~8{,}300). The Fiedler eigenvalue declines
monotonically from 0.0019 to 0.0004, entering the same regime as the
aggressive-$c_w$ experiments. The growth rate \emph{accelerates} with
accumulated BD strength, suggesting a positive feedback mechanism: once
block-diagonal structure exceeds a critical threshold (${\sim}2{,}000$:1),
the contrastive loss becomes more effective at reinforcing the existing
partition, creating a self-amplifying cycle.

\textbf{Phase~4 (steps 8{,}300--10{,}000): BD regime stabilization.}
The co/cross ratio oscillates between 12{,}145:1 and 15{,}183:1
(mean~13{,}288:1) without a clear growth trend. The Fiedler eigenvalue
stabilizes at 0.00042--0.00048. Validation loss continues declining
(0.4034 $\to$ 0.4017), indicating that the \emph{language model} is still
improving even as the \emph{bridge structure} has reached equilibrium.
This decoupling---structural stasis while perplexity improves---suggests
the bridge has settled into an attractor from which the remaining
training budget refines weights within the established topology rather
than reshaping it.

\subsubsection{Speed, Destination, or Both?}

At step~10{,}000, CW-001's Fiedler eigenvalue (0.00046) is within one
order of magnitude of H-ch6's final Fiedler (0.00009). Both are deep
in the BD regime. But the co/cross ratios tell a more nuanced story:

\begin{table}[h]
\centering
\small
\begin{tabular}{@{}lrrr@{}}
\toprule
Metric & H-ch6 ($c_w = 0.1$ adaptive) & CW-001 ($c_w = 0.005$ floor) & Ratio \\
\midrule
Co/cross & 70{,}404:1 & 13{,}456:1 & $5.2\times$ \\
Fiedler & 0.00009 & 0.00046 & $5.1\times$ \\
Val loss & ${\sim}0.40$ & 0.4017 & ${\sim}$same \\
\bottomrule
\end{tabular}
\caption{CW-001 vs.\ H-ch6 at matched step count (10{,}000 steps).}
\label{tab:cw001-comparison}
\end{table}

A $5.2\times$ gap persists at matched step count. CW-001's Phase~4
stabilization---oscillating at 12{,}145--15{,}183:1 for 2{,}000~steps
with no growth trend---raises the question: is this a transient plateau
(like Phase~2, which lasted 2{,}700~steps before breaking out) or a
genuine attractor at this $c_w$ floor?

Two interpretations compete. Under the \textbf{speed hypothesis}, CW-001
simply needs more training steps; the Phase~$2 \to 3$ breakout precedent
shows the system can appear stalled and then resume exponential growth.
Under the \textbf{attractor hypothesis}, the contrastive floor sets not
only the rate but the ceiling of structural formation. The $5.2\times$
co/cross gap and the $5.1\times$ Fiedler gap are proportional, consistent
with a $c_w$-dependent attractor strength.

The 24C-002 experiment (\S\ref{sec:24cell}) resolves this question in
favour of the attractor hypothesis. At $n{=}24$ with adaptive $c_w$
settling at $0.025$, the final co/cross is $5{,}983$:1---$6\times$ below
the fixed-$c_w{=}0.1$ run (24C-001, $35{,}808$:1) at matched step count.
The contrastive weight governs not only speed but the \emph{depth} of BD
formation. Higher $c_w$ produces a deeper attractor.

Both interpretations share the same novel core: \emph{any nonzero $c_w$
produces BD structure}, the breakout mechanism at ${\sim}2{,}000$:1 is a
genuine phase transition, and the three-phase incubation trajectory
has no prior art in the adapter literature.

\subsubsection{The Phase~2 Plateau as Incubation}

The 3{,}000-step plateau (Phase~2) is not a failure state. During this
period, the Fiedler eigenvalue rises from 0.0004 to 0.0020---the bridge
is building spectral connectivity while block-diagonal strength appears
stalled. We hypothesize that Phase~2 represents an \emph{incubation
period} in which the bridge accumulates the spectral prerequisites for
exponential BD growth. The Fiedler rebound creates a connected substrate;
Phase~3's breakout then partitions this substrate into blocks.

This two-stage process (connect, then partition) mirrors the
spectral-attractor-then-bifurcation sequence observed across the full
experimental programme (\S\ref{sec:spectral-attractor}--\ref{sec:negative-controls}):
spectral-only training creates connectivity (the attractor); contrastive
loss then breaks this connectivity into directed blocks. CW-001
reproduces this sequence \emph{within a single run} at low~$c_w$.

\subsubsection{Implications}

\begin{enumerate}
\item \textbf{No minimum $c_w$ threshold for BD.} Any nonzero $c_w$
  eventually produces BD structure---the FI-004 ``cliff'' at $c_w = 0$
  is real, but $c_w = 0.005$ is far above it.
\item \textbf{Adaptive decay is not a bottleneck.} The Steersman's
  adaptive Control Law~2, which reduces $c_w$ when BD is detected, does
  not prevent BD formation---it merely slows it. The floor mechanism
  ensures a minimum contrastive signal is always present.
\item \textbf{Training budget trades against $c_w$.} A practitioner with
  abundant compute can use lower $c_w$ for potentially better spectral
  properties; a practitioner needing fast convergence can use higher
  $c_w$. Whether the slow regime produces qualitatively better adapters
  (lower Fiedler $=$ more connected $=$ better information flow) is an
  open empirical question.
\item \textbf{The breakout threshold (${\sim}2{,}000$:1) is a phase
  transition.} The co/cross ratio of ${\sim}2{,}000$:1 marks the onset
  of self-amplifying BD growth at $c_w = 0.005$. Pending verification
  across different $n$ and topologies, this may be universal.
\item \textbf{BD regime stabilization decouples from language modeling.}
  Validation loss continues declining (0.4034 $\to$ 0.4017) during
  Phase~4 while co/cross oscillates without growth. The bridge structure
  reaches equilibrium before the language model does---the final training
  budget refines representations within established topology.
\end{enumerate}
